\subsection{\texorpdfstring{\(\mathcal{L}(\mathbf{E};\mathbf{F})\) 的等度连续子集}{\textbackslash mathcal\{L\}(\textbackslash mathbf\{E\};\textbackslash mathbf\{F\}) 的等度连续子集}}

设 \(\mathrm{E}\) 和 \(\mathrm{F}\) 是两个局部凸空间。为使子集 \(\mathrm{H}\)（\(\mathcal{L}(\mathrm{E};\mathrm{F})\) 的）等度连续，必要且充分的是它在 \(\mathrm{E}\) 中的点 0 处等度连续（I, 第 9 页, 命题 6）；这蕴含：对 0 的每个邻域 \(\mathrm{V}\)（在 \(\mathrm{F}\) 中），集 \(\bigcap_{u\in \mathrm{H}}u^{-1}(\mathrm{V})\) 是 \(\mathrm{E}\) 中 0 的一个邻域；或者，对每个连续半范数

\(p\)（\(\mathrm{F}\) 上的），函数 \(\sup_{u\in \mathrm{H}}(p\circ u)\) 是 E 上的连续半范数。此外（I, 第 5 页），

  H 是一致等度连续的。我们注意到，等度连续子集的凸均衡包是等度连续的，因为若 p 是 F 上的连续半范数而 \(\tilde{H}\) 是 H 的凸均衡包，则对 H 中的 \(u_{i}\)，有不等式 \(p \circ (\sum_{i} \lambda_{i} u_{i}) \leqslant \sum_{i} |\lambda_{i}| \cdot (p \circ u_{i})\)，因而 \(\sup_{u \in \tilde{H}} (p \circ u) = \sup_{u \in H} (p \circ u)\)。

因此，等度连续子集的族是 \(\mathcal{L}(\mathrm{E};\mathrm{F})\) 上的一个凸有界结构（III, 第 2 页, 定义 2）。

命题 4. --- 设 E、F 是两个局部凸空间，且 F 是 Hausdorff 的。给从 E 到 F 中的所有映射所成的空间 \(F^{E}\) 赋予简单收敛拓扑。则

\begin{enumerate}
\def\labelenumi{(\roman{enumi})}
\item
  从 \(\mathbf{E}\) 到 \(\mathbf{F}\) 中的线性映射所成之集在 \(\mathbf{F}^{\mathrm{E}}\) 中是闭的。
\item
  若 \(\mathbf{H}\) 是 \(\mathcal{L}(\mathrm{E};\mathrm{F})\) 的一个等度连续子集，则闭包 \(\overline{\mathbf{H}}\)（\(\mathbf{H}\) 在 \(\mathrm{F}^{\mathrm{E}}\) 中的）含于 \(\mathcal{L}(\mathrm{E};\mathrm{F})\) 且是等度连续的。
\end{enumerate}

我们知道 \(\overline{\mathbf{H}}\) 是等度连续的（GT, X, § 2, No.~3, 命题 6）。剩下要证明断言 (i)。设 \(x, y\) 在 E 中且 \(\lambda, \mu\) 在 K 中，并设 \(\mathbf{A}(x, y, \lambda, \mu)\) 是所有满足下式的 \(u \in \mathbf{F}^{\mathrm{E}}\) 所成之集：

\[
u (\lambda x + \mu y) - \lambda u (x) - \mu u (y) = 0.
\]

这个集在 \(F^{E}\) 中是闭的，因为映射 \(u \mapsto u(x)\)（从 \(F^{E}\) 到 F 中）对每个 \(x \in E\) 连续，且 \(F\) 是 Hausdorff 的。但从 \(E\) 到 \(F\) 中的线性映射所成之集等于

\[
\bigcap_ {x, y, \lambda , \mu} \mathrm{A} (x, y, \lambda , \mu).
\]

于是这个集在 \(F^{E}\) 中是闭的。

推论 1. --- 为使 \(\mathcal{L}(\mathrm{E};\mathrm{F})\) 的等度连续子集 H 在 \(\mathcal{L}_s(\mathrm{E};\mathrm{F})\) 中相对紧，必要且充分的是：对所有 \(x\in \mathrm{E}\)，集 \(\mathrm{H}(x)\)（由所有 \(u(x)\)（当 \(u\) 取遍 \(\mathrm{H}\)）组成）在 \(\mathrm{F}\) 中相对紧。

事实上，这个条件对 \(\overline{\mathbf{H}}\) 在 \(\mathrm{F}^{\mathrm{E}}\) 中紧是必要且充分的（GT, I, § 9, No.~5, 推论）。

推论 2. --- 对偶 E' 的每个等度连续子集在 E' 上的弱拓扑 \(\sigma(E', E)\) 下相对紧（III, 第 14 页, 例 4）。

因为若 H 是 \(E'\) 的一个等度连续子集，则 \(\sup_{u\in H}|u|\) 是 E 上的连续半范数；特别地，对每个 \(x\in E\)，集 \(\mathrm{H}(x)\) 有界，从而在标量域中相对紧。

推论 3. --- 在半赋范空间 E 的强对偶 \(E_{b}^{\prime}\) 中，每个闭球对弱拓扑 \(\sigma(\mathrm{E}^{\prime},\mathrm{E})\) 都是紧的。

这个球对 \(\sigma(\mathrm{E}^{\prime},\mathrm{E})\) 也是闭的。

命题 5. --- 设 E 和 F 是两个局部凸空间，并设 T 是 E 的一个完全子集。下列一致结构在 \(\mathcal{L}(\mathrm{E};\mathrm{F})\) 的每个等度连续子集 H 上重合：

\begin{enumerate}
\def\labelenumi{\arabic{enumi})}
\item
  T 上的简单收敛一致结构；
\item
  E 上的简单收敛一致结构；
\item
  在 E 的预紧子集上收敛的一致结构。
\end{enumerate}

我们回忆（III, 第 15 页, 命题 2）：\(\mathfrak{S}\) -拓扑（\(\mathcal{L}(\mathrm{E};\mathrm{F})\) 上的）与 \(\tilde{\mathfrak{S}}\) -拓扑重合，其中 \(\tilde{\mathfrak{S}}\) 是适配于 E 且包含 \(\mathfrak{S}\) 的最小有界结构。

因此在命题 5 的陈述中，我们可以把「完全」一词换成「处处稠密」。于是该命题由等度连续集的一般性质得出（GT, X, § 2, No.~4, 定理 1）。

例. --- * 1) 设 \(\mu\) 是 \(\mathbf{R}\) 上的 Lebesgue 测度，并设 E 是半赋范空间 \(\mathcal{L}^p(\mu)\)（\(1 \leqslant p < \infty\)）（INT, IV）。对每个数值函数 \(f\) 和每个实数 \(h\)，设 \(f_h\) 为函数 \(x \mapsto f(x - h)\)。显然映射 \(f \mapsto f_h\) 定义了 E 到它自身上的一个线性等距。若 \(f\) 连续且有紧支集，则 \(f_h\) 一致收敛于 \(f\)，从而也以 \(p\) 阶平均收敛（当 \(h\) 趋于 0 时）。由于所有具有紧支集的连续函数所成之集 \(\mathcal{K}(\mathbf{R})\) 在 E 中稠密，且 E 的线性等距所成之集是等度连续的，由命题 5 可得：对每个 \(f \in E\)，\(f_h\) 以 \(p\) 阶平均收敛于 \(f\)（当 \(h\) 趋于 0 时）。

对 \(p = 1\)，考虑 Fourier 变换，它使每个 \(f \in \mathcal{L}^1(\mu)\) 对应于函数 \(\hat{f}\)（定义在 \(\mathbf{R}\) 上），后者由下式定义

\[
\hat {f} (y) = \int e ^ {- 2 i \pi x y} f (x) d \mu (x).
\]

线性型 \(f \mapsto \hat{f}(y)\) 所成之集是 \(\mathcal{L}^1(\mu)\) 的对偶的一个等度连续子集。另一方面，我们知道闭有界区间的特征函数所成之集 \(T\) 是 \(\mathcal{L}^1(\mu)\) 的一个完全子集；并且容易验证：对所有 \(f \in T\)，Fourier 变换 \(\hat{f}\) 是一个连续函数且在无穷远处趋于零。我们推出这对所有 \(f \in \mathcal{L}^1(\mu)\) 都成立（「Riemann-Lebesgue 定理」）。

关系式 \(\sup_{y\in \mathbb{R}}|\hat{f} (y)|\leqslant \| f\| _1\) 表明映射 \(f\mapsto \hat{f}\) 是一个连续映射

它从 \(\mathcal{L}^1 (\mu)\) 映入空间 \(\mathcal{B}(\mathbf{R})\)（即 \(\mathbf{R}\) 上所有有界函数所成的空间），后者带有一致收敛结构。由于 \(\hat{f}\) 对所有 \(f\in \mathrm{T}\) 连续，由此得出 \(\hat{f}\) 对所有 \(f\in \mathcal{L}^1 (\mu)\) 连续。\(\hat{f}\) 在无穷远处趋于零这一事实由如下事实得出：在无穷远处趋于零的所有连续函数所成的子空间 \(\mathcal{C}_0(\mathbf{R})\) 在 \(\mathcal{B}(\mathbf{R})\) 中是闭的。

\begin{enumerate}
\def\labelenumi{\arabic{enumi})}
\setcounter{enumi}{1}
\tightlist
\item
  设 E 是 R 上所有连续数值函数所成的空间，带上紧收敛拓扑。设 K 是 R 的一个紧子集，并设 \((\mu_{n})\) 是 R 上支集含于 K 的测度的一个序列。假设对所有 n 有 \(\|\mu_{n}\|\leqslant1\)。于是诸 \(\mu_{n}\) 所成之集是 \(E'\) 的一个等度连续子集。因此，若对每个函数 \(f\in E\) 有 \(\lim_{n\to\infty}\mu_{n}(f)=0\)，则函数序列 \(x\mapsto\int e^{itx}d\mu_{n}(t)\) 收敛于 0，
\end{enumerate}

且在 \(\mathbf{R}\) 的每个紧子集上一致收敛（因为诸函数 \(t \mapsto e^{itx}\)——当 \(x\) 取遍 \(\mathbf{R}\) 的一个紧子集时——所成之集在 E 中是紧的）。*

推论. --- 假设 F 是 Hausdorff 的。设 H 是 \(\mathcal{L}(\mathrm{E};\mathrm{F})\) 的一个等度连续子集。若 H 上的一个滤子 \(\Phi\) 简单收敛到一个从 E 到 F 中的映射 \(u_{0}\)，则 \(u_{0}\) 是从 E 到 F 中的一个连续线性映射，且 \(\Phi\) 在 E 的每个预紧子集上一致收敛于 \(u_{0}\)。

第一个断言由命题 4（III, 第 16 页）得出，第二个由命题 5（III, 第 17 页）得出。

命题 6. --- 设 H 是 \(\mathcal{L}(\mathrm{E};\mathrm{F})\) 的一个等度连续子集。若 F 可度量化且 E 中存在一个可数完全集，则 H 上 E 中简单收敛的一致结构是可度量化的。若此外 F 中还存在一个可数完全集，则 H 中存在一个可数的处处稠密集（对 E 的紧子集上的一致收敛拓扑而言）。

设 \((a_{n})\) 是 E 中的一个完全序列。则映射 \(u \mapsto (u(a_{n}))\) 是从 \(\mathcal{L}(\mathrm{E};\mathrm{F})\)（其中 \(\mathcal{L}(\mathrm{E};\mathrm{F})\) 带有诸 \(a_{n}\) 所成之集上简单收敛的一致结构）到 \(F^{N}\) 的一个一致子空间上的一个同构。若 F 可度量化（相应地，可度量化且满足第一可数公理），则 \(F^{N}\) 亦然（GT, IX, § 2, No.~4, 推论 2 及 § 2, No.~8, 推论），而命题由命题 5（III, 第 17 页）得出。

推论 1. --- 设 E 是一个局部凸可度量化空间，F 是一个赋范空间。假设 E 和 F 都满足第一可数公理。则 \(\mathcal{L}(\mathrm{E};\mathrm{F})\) 是可数个等度连续子集所成的族的并，且在 \(\mathcal{L}(\mathrm{E};\mathrm{F})\) 中存在一个可数集，它对 E 的预紧子集上的一致收敛拓扑是稠密的。

设 B 是 F 的单位球，\((\mathrm{V}_{n})\) 是 E 中 0 的一个可数基本邻域系。对每个整数 n，集 \(H_{n}\)（由所有 \(u \in \mathcal{L}(\mathrm{E}; \mathrm{F})\) 中使 \(u(\mathrm{V}_{n}) \subset \mathrm{B}\) 者组成）是等度连续的，且 \(\mathcal{L}(\mathrm{E}; \mathrm{F})\) 是诸 \(H_{n}\) 的并。于是该推论由命题 6 得出。

推论 2. --- 在满足第一可数公理的赋范空间的对偶 E' 中，每个闭球对弱拓扑 \(\sigma(\mathrm{E}',\mathrm{E})\) 都是一个紧可度量化空间，且对该拓扑，E' 中存在一个可数稠密子集。

这由命题 6 以及 III, 第 17 页, 推论 3 得出。

